\documentclass[11pt, a4paper]{article}
\usepackage[utf8]{inputenc}
\usepackage[T1]{fontenc}
\usepackage[french]{babel}
\usepackage{graphicx}
\usepackage[normalem]{ulem}
\usepackage{amsmath}
\usepackage{amssymb}
\usepackage{xcolor}
\usepackage[backend=biber]{biblatex}

\author{}
\date{}
\title{}

\begin{document}

%% Contexte

Dans Cascaval (2023), les auteurs proposent de répondre au problème du
référencement d'éléments géométriques dans les opérations de modélisation
et placent leurs travaux dans un contexte défini par les concepts suivants :
\begin{itemize}
	\item Un modèle paramétrique est une fonction prenant des paramètres
	      géométriques et scalaires.
	\item Une référence peut être obtenue grâce à une requête sur un élément
	      géométrique.
	\item Les types géométriques traités sont les point, segment (ligne ou
	      arc), polygone (conceptuellement une face sans trou) et face (un
	      élément produit à partir d'un ou plusieurs polygones).
\end{itemize}

Dans ce contexte, un langage conçu pour la CAO doit respecter trois propriétés :
\begin{itemize}
	\item il ne doit pas introduire de discontinuité dans le référencement d'une
	      géométrie,
	\item il doit permettre la sélection de tout ou partie d'un ensemble d'éléments
	      géométrique résultant de la partition d'une géométrie passée en entrée d'une
	      opération,
	\item l'utilisateur doit pouvoir exclure des configurations topologiques des modèles
	      en sortie du programme.
\end{itemize}

%% Contribution

Dans ce contexte, Cascaval et al. proposent un langage spécifique à la
conception de modèles paramétriques mettant en œuvre un système de généalogie
avec une sémantique permettant de décrire avec une certaine précision les
relations entre divers éléments géométrique respectivement en entrées et en
sortie d'une opération. Dans ce langage, une référence est définie par un
utilisateur qui crée une requête sur un élément géométrique. Ces requêtes
peuvent présenter, de manière explicite, des contraintes liées à cette
généalogie pour représenter précisément les intentions de modélisation de
l'utilisateur.

%% Explication de la sémantique

À cette fin, trois relations sont définies :
\begin{itemize}
	\item la relation de transformation, de sous-ensemble et de dérivation. La
	      relation de transformation établie un lien directe entre un élément en
	      entrée et en sortie à travers une opération. En effet, les opérations
	      telles que la rotation et la translation ne créent rien de nouveau que
	      ce soit en termes de géométrie ou de topologie. Seuls les paramètres
	      géométriques (grandeur, position, …) sont modifiés.
	      {\color{red}accompagner d'une illustration pour servir d'exemple}
	\item La relation de sous-ensemble établie un lien entre un ensemble
	      d'éléments en entrée d'une opération et un élément en sortie. Il s'agit,
	      typiquement dans les opérations de CSG, de faire le lien entre, par
	      exemple deux arêtes qui s'intersectent et un sommet résultat de cette
	      intersection.
		      {\color{red}accompagner d'une illustration pour servir d'exemple}
	\item La relation de dérivation, elle, établie plusieurs liens entre de la
	      géométrie existante et de la géométrie nouvellement créée. Prenons le
	      cas du chanfrein sur un sommet présenté par les auteurs. Lorsque
	      l'opération de chanfrein est appliquée, une arête est créée à partir
	      d'un sommet, il s'agit du premier lien établi. Il faut ensuite faire la
	      distinction entre les deux sommets aux extrémités de cette arête. Pour
	      cela, la relation de dérivation établit un lien entre l'un des deux
	      sommets une arête adjacente pré-existante.
	      {\color{red}accompagner d'une illustration pour servir d'exemple}
\end{itemize}

%% Principe et fonctionnement des requêtes

Une requête, un ensemble de clauses logiques avec des notions d'ensembles,
permet de définir une référence à partir d'un élément géométrique. %
On peut, par exemple, récupérer tout (\texttt{all}) ou partie (\texttt{from(e)})
d'un élément e, de tous les éléments du modèle (\texttt{fromAll(e)}), d'un
élément parmi un ou plusieurs ensembles (\texttt{fromAny((e | \{e\})*)}). Il est
également possible de créer des conjonctions de requêtes q (\texttt{and(q,q)}),
des disjonctions (\texttt{or{q,q}}) et même en inclure (\texttt{contains(q)}) et
en exclure (\texttt{not(q)}). %
L'utilisateur peut s'appuyer sur la sémantique de généalogie pour décrire au
mieux ses intentions de modélisation. À ce moment là, les requêtes expriment,
explicitement, les relations de généalogie décrites plus tôt. Par exemple, les
éléments dérivés d'une géométrie lors d'une opération peuvent être obtenues en
remplaçant \texttt{from} par \texttt{derivedFrom} dans toutes les clauses
d'ensembles (\texttt{from, fromAll, fromAny} deviennent \texttt{derivedFrom,
	derivedFromAll, derivedFromAny}). %
	{\color{red}accompagner d'illustrations pour servir d'exemples}

%%% Continuité des références et gestion de l'ambiguïté

Dans les cas de l'ambiguïté est créée, par exemple, de l'intersection de
segments, des propriétés géométriques peuvent être ajoutées aux éléments
géométriques pour lever cette ambiguïté. %
Dans le cas de l'intersection de lignes, chacun des points créés doit être
considéré comme l'extrémité \texttt{end} de l'une des deux lignes. %
Pour l’intersection d'une ligne avec un arc, la distance entre les points de la
ligne est normalisée, et les points d'intersections sont enrichis avec leur
distance au sommet de départ (\texttt{start}) de la ligne. %
Dans le cas de l’intersection de deux arcs, Cascaval et al. reprennent les
travaux de Bidarra. Une ligne traversant les deux arcs en leur centre crée un
domaine positif et un domaine négatif pour enrichir les intersections, notons
que ces domaines ne varient pas, peu importe les transformations sur le modèle. %

%% Comparaison avec Cascaval

Ici, Cascaval et al. proposent un langage spécifique pour la conception de
modèles paramétriques. Ce langage présente diverses fonctionnalités telles que
la définition de références géométrique à partir de requêtes utilisateurs, les
éléments géométriques étant enrichis d'information géométrique et généalogique
pour garantir le respect des intentions de modélisation de l'utilisateur.
Néanmoins, cette approche a plusieurs inconvénients :
\begin{itemize}
	\item Les éléments géométriques utilisés pour les références sont des cellules
	      dont le traitement n'est pas généralisable en toutes dimensions.
	\item Pour la sémantique pour la généalogie, bien que les auteurs aient mis en
	      place un système pour tracer une forme d'origine pour distinguer une
	      géométrie d'une autre, cela ne leur permet en fin de compte de ne
	      distinguer que la création, la suppression et la modification au sens
	      large de la topologie d'une géométrie.
	\item Les opérations de modélisation sont responsables de l'annotation de la
	      géométrie pour le suivi. Ce qui implique que l'implantation du code de
	      ces opérations peut être source d'erreur dans le suivi.
	\item L'utilisateur est contraint de définir lui-même les requêtes de
	      définition d'une référence, cela peut compliquer l'automatisation du
	      processus de réévaluation notamment lorsqu'il y a ajout, suppression,
	      déplacement d'opération.
\end{itemize}


\end{document}
